% !TEX root = chapter-standalone.tex

\section{Monotone relations and design problems}

\todotext{J: introduce DPs already here}


\begin{definition}[monotone relation]
Let $\posAdefinition$ and~$\posBdefinition$ be posets. A monotone relation $\relA \colon \posA \mto \posB$ is a relation $\relA \colon \posAset \to \posBset$ such that for all $\ela, \ela' \in \posAset$ and all $\elb, \elb' \in \posBset$, 
\begin{enumerate}
\item $\tup{\ela, \elb} \in \relA, \ \ela' \leq_{\posA} \ela \ \implies \tup{\ela', \elb} \in \relA$;
\item $\tup{\ela, \elb} \in \relA, \ \elb \leq_{\posB} \elb' \ \implies \tup{\ela, \elb'} \in \relA$.
\end{enumerate}
\end{definition}

\subsection{Monotone relations and co-design}

A monotone relation
\begin{equation}
\adp \colon \funsp \mto \ressp
\end{equation}
can be used to model a relationship of ``feasability'' between a poset $\funsp$ of ``functionalities'' and a poset $\ressp$ of ``requirements", in the sense that the relation describes whether a resource $\fun \setin \funsp$, seen as a functionality or service or product, is \emph{feasible} to obtain given a certain resource $\res \setin \ressp$, with $\res$ interpreted as a requirement or a cost.

The condition
\begin{equation}\label{eq:mon-rel-cond-1}
\tup{\fun, \res} \in \adp, \ \res \leq_{\ressp} \res' \ \implies \tup{\fun, \res'} \in \adp
\end{equation}
says that if $\fun$ is feasible to obtain using $\res$, then it is also feasible to obtain $\fun$ if we use more resources, $\res'$. 

The condition
\begin{equation}\label{eq:mon-rel-cond-2}
\tup{\fun, \res} \in \adp, \ \fun' \leq_{\funsp} \fun \ \implies \tup{\fun', \res} \in \adp,
\end{equation}
on the other hand, says that if $\fun$ is feasible to obtain using $\res$ amount of resources, then it is also feasible to obtain less, $\fun'$, using the same resources $\res$. 



\subsection[Design problems]{Design problems}
\label{sec:dpdefinition}
\SYNDEF{design problem}

Here is an alternative way to think about relations of feasibility. Given a particular functionality~$\fun \in \funsp$ and requirement~$\res \in \ressp$, we would like to know a ``true'' or ``false'' answer to the question of whether they form a feasible pair of resources. This situation is described by a function
%
\begin{equation}
    \mapb\colon \funsp \cartprod \ressp \sto \boolset.
    % leave cartprod
\end{equation}
%
The value~$\mapb(\fun, \res)$ is the answer to the question ``is the functionality~\fun feasible with resources~\res?''.

The conditions \cref{eq:mon-rel-cond-1} and \cref{eq:mon-rel-cond-2} can now be translated into this formulation: they say, respectively, that $\mapb\colon \funsp \cartprod \ressp \sto \boolset$ is monotone in the variable $\res$ and antitone in the variable $\fun$. Or, equivalently, we can say that we have a monotone map of two variables of the type
\begin{equation}
     \funspop \Ptimes \ressp \sto \Bool.
    % leave cartprod
\end{equation}



%\linkvideo{spring2021-design:design:bool-prof} % Boolean profunctors

\begin{definition}[Design Problem]
    \label{def:design-problem}
    Given posets $\funsp$ and $\ressp$, a \maindef{design problem} (DP) from $\funsp$ to $\ressp$ is a \SY{monotone map} of the form%
    \begin{equation}
        \label{eq:design_prob_as_monfun}
        \adp \colon \funspop \Ptimes \ressp \toinPos \Bool.
    \end{equation}
\end{definition}

\begin{definition}[Feasible set of a design problem]
    \label{def:dp-feasible-set}
    We define the \emph{feasible set}~$\feasibleset{\adp}$ of a design problem
    \begin{equation}
        \label{eq:generic-design-prob}
        \adp\colon \funspop \Ptimes \ressp \toinPos \Bool
    \end{equation}
    as the subset of~$\funspop\Ptimes \ressp$ for which~$\adp$ is the \emph{indicator function}, that is%
    \begin{equation}
        \label{eq:design-prob-as-upper-set}
        \feasibleset{\adp} = \makeset{ \tup{\F{\fun^*},\res} \setin \funspop\Ptimes \ressp \mid %
            \adp(\fun\Fop, \res) = \true %
        }.
    \end{equation}
\end{definition}

Note that the feasibility set~$\feasibleset{\adp}$ of a \SY{design problem}~$\adp \colon \funspop \Ptimes \ressp\toinPos \Bool$ is a binary relation~$\feasibleset{\adp} \setsubseteq \funspop \Ptimes \ressp$.
We saw in \cref{rem:rel-three-fun-descriptions} that there is a one-to-one correspondence between functions~$\mapb \colon \setA \cartprod \setB \sto \boolset$ and binary relations~$\relA\colon \setA \mto \setB$.




\begin{gradedexercise}[\exname{DPsAsUpperSets}]
    \label{ex:DPsAsUpperSets}

    In this exercise, your task is to prove that there is a one-to-one correspondence between \SY{design problems}~$\adp \colon \funspop \Ptimes \ressp\toinPos \Bool$ and \SY{upper sets}~$\feasibleset{} \setsubseteq \funspop \Ptimes \ressp$.

    In more detail:

    Let $\setA$ denote the set of all \SY{design problems}~$\adp \colon \funspop \Ptimes \ressp\toinPos \Bool$ and let $\setB$ denote that set of all \SY{upper sets}~$\feasibleset{} \setsubseteq \funspop \Ptimes \ressp$.

    \begin{enumerate}
        \item Define a function
              \begin{equation*}
                  \mapa \colon \setA \sto \setB
              \end{equation*}
              which to any design problem in $\setA$ assigns a corresponding upper set in $\setB$.

        \item Define a function
              \begin{equation*}
                  \mapb \colon \setB \sto \setA
              \end{equation*}
              which maps any upper set in $\setB$ to a corresponding design problem in $\setA$.
        \item Prove that $\mapa$ and $\mapb$ are inverses to one another.
    \end{enumerate}
\end{gradedexercise}

\solutionof{DPsAsUpperSets}


\begin{gradedexercise}[\exname{DPsFromMonotoneMaps}]
    \label{ex:DPsFromMonotoneMaps}

    Given any monotone map~$\mapb\colon \funsp \toinPos \ressp$, we can turn it into a design problem
    \begin{equation*}
        \adp_\mapb \colon \funspop \Ptimes \ressp \toinPos \Bool
    \end{equation*}
    via the following recipe.
    Set
    \begin{equation*}
        \adp_\mapb (\fun\Fop, \res) = \true \quad \text{if and only if} \quad \mapb(\fun) \posleq \res.
    \end{equation*}

    Prove that~$\adp_\mapb$, as defined above, is indeed a \SY{design problem} when~$\mapb$ is a \SY{monotone map}.
\end{gradedexercise}

\solutionof{DPsFromMonotoneMaps}

Recall that when working with a relation~$\relA \colon \setA \sto \setB$ between sets, if the sets in question were finite, then we could conveniently draw the relation~\relA using arrows to connect elements of~\setA to those elements of~\setB to which they are related via~\relA.
Given a \SY{design problem}~$\adp \colon \funspop \Ptimes \ressp\toinPos \Bool$ involving finite \SY{posets}, we can visualize it in a similar fashion.
We use Hasse diagrams to visualize the \SY{posets} involved, and we use dashed arrows to connect those elements which are related via the feasibility set~$\feasibleset{\adp}$.

\begin{marginfigure}
    \centering
    \includesag{example_dp_composition_xy}
    \caption{}
    \label{fig:example_dp_graph_xy}
\end{marginfigure}
\begin{marginfigure}
    \centering
    \includesag{example_dp_graph_xy_feas}
    \caption{}
    \label{fig:example_dp_graph_xy_feas}
\end{marginfigure}

\begin{example}
    \label{exa:visualize-dp}
    In \cref{fig:example_dp_graph_xy} we have illustrated this kind of visualization in the case of a \SY{design problem} of the type
    \begin{equation}
        \adp \colon \F{\posgenA} \posop \Ptimes \R{\posgenB} \toinPos \Bool,
    \end{equation}
    where~$\F{\posgenA} = \tup{\posAset, \posleq}$ and~$ \R{\posgenB} = \tup{\posBset, \posleq}$ are finite \SY{posets}, with
    \begin{equation}
        \posAset = \makeset{ \setAeln{1}, \setAeln{2}, \setAeln{3} }
        \qqand
        \posBset = \makeset{ \setBeln{1}, \setBeln{2}},
    \end{equation}
    and ordered as shown in the figure.

    The relation described by the \SY{design problem} is marked with the dashed arrows;
    The feasibility set
    \begin{equation}
        \feasibleset{\adp} = \makeset{\tup{\setAeln{1}, \setBeln{1}}, \tup{\setAeln{1}, \setBeln{2}}, \tup{\setAeln{2}, \setBeln{2}}, \tup{\setAeln{3}, \setBeln{2}} },
    \end{equation}
    is reported in \cref{fig:example_dp_graph_xy_feas}.
\end{example}




The Boolean-valued \SY{design problems} we are considering in this section do not distinguish between particular implementations: they only tell us if \emph{any} implementation or solution exists for given functionality and resources.


    \paragraph{Diagrammatic notation}
    We represent \SY{design problems} using a diagrammatic notation.
    One \SY{design problem}~$\adp \colon \funsp \profto \ressp$ is represented as a box with functionality~\funsp on the \emph{left} and resources~\ressp on the \emph{right} (\cref{fig:diagrammaticdp}).
    \begin{figure}[h!]
        \centering
        \includesag{50_diagrammatic}
        \caption{Diagrammatic representation of a design problem. }
        \label{fig:diagrammaticdp}
    \end{figure}

    As we did for DPIs, we will connect these diagrams.


\begin{marginfigure}
    \centering
    \includesag{50_engine}
    \caption{Diagram of the engine design problem.}
    \label{fig:enginedp}
\end{marginfigure}

    \begin{example}
        An aerospace company, Jeb's Spaceship Parts, is designing a new rocket engine, the Bucket of Boom X100.
        The engine requires fuel and provides thrust, and so it can be modeled as a \SY{design problem} where \R{fuel} and \F{thrust} are two totally-ordered sets representing their respective resources.

        The corresponding diagram is reported in~\cref{fig:enginedp}.

        Concretely, ``engine'' is represented as a monotone map%
        %
        \begin{equation}
            \text{engine} \colon \Ftext{thrust}\posop \Ptimes \Rtext{fuel} \toinPos \Bool.
        \end{equation}
        %
        Assuming that the \SY{posets} \R{fuel}, \F{thrust}$\op$ are finite, we can think of the ``engine'' \SY{design problem} as a matrix, where each~$(i,j)$-th entry is the answer to the question, ``is the amount of thrust~$\F{f_i}$ feasible with the amount of fuel~$\R{r_j}$?'':
        %
        \begin{equation}
            \begin{blockarray}{ccccccc}
                &&&& \Rtext{Fuel} \\
                && \R{r_1} = 0 & \R{r_2} & \R{r_3} & \hdots & \R{r_m} \\
                \begin{block}{cc(ccccc)}
                    & \F{f_n}\Fop = 0 & 0 & 0 & 0 & & 0 \\
                    & \F{f_{n-1}^*} & 0 & 0 & 0 & & 1 \\
                    \Ftext{thrust}\op & \F{f_{n-2}^*} & 0 & 1 & 1 & & 1 \\
                    & \vdots & & & & \ddots & \\
                    & \F{f_1^*} & 1 & 1 & 1 & & 1 \\
                \end{block}
            \end{blockarray}
        \end{equation}
        %
        Suppose we have tested or are given the performance data of a few different engines, as possible solutions to the ``engine'' design problem, each with a fixed optimal fuel-thrust value.
        To illustrate the monotonicity assumption, we can render the data of ``engine'' as a graph, as depicted in~\cref{fig:solenginedp}.
        \begin{figure}[h!]
            \centering
            \includesag{50_engine_graph}
            \caption{Graphical representation of the possible solutions of the engine design problem. }
            \label{fig:solenginedp}
        \end{figure}

        Note that the shaded regions cover the feasible solution set.
        This feasible solution set is always an \emph{upper set} (\cref{def:upperset}) in~$\Ftext{thrust}\posop \Ptimes \Rtext{fuel}$, which is another way of characterizing the monotonicity of the design problem.
        The optimal solutions, indicated by dots, form an \emph{antichain} of solutions.
        We will come back to \SY{antichains} when discussing how to compute optimal solutions of \SY{design problems}.

    \end{example}


    

%\JT{Do we need 2 large examples, or can we move the battery example to the notes?}
%
%\newcommand{\cCapacity}{\text{Capacity}}
%\newcommand{\cMass}{\text{Mass}}
%
%\begin{example}[Battery]
%   A battery is a store of elecrical energy.
%   We are interested in the relation between the capacity of the battery, measured
%   in joules, and the mass of the battery, measured in grams.
%   We will model the battery as a design problem
%   \[
%       \text{battery} : \cCapacity \profto \cMass,
%   \]
%   with the two \SY{posets} $\cCapacity$ and $\cMass$ defined as
%   $\cCapacity \definedas \langle\mathbb{R}_+^{\text{J}}, \leq\rangle$
%   and $\cMass \definedas \langle\mathbb{R}_+^{\text{g}}, \leq\rangle$, where ``$\leq$''' is the usual
%   order on $\mathbb{R}_+$.
%
%This is the corresponding diagram:
%   \[
%   \centering
%   \begin{tikzpicture}[oriented WD, bb min width =1.5cm, bby=2ex, bbx=.7cm,bb port length=3pt]
%       \node[bb port sep=0.8, bb={1}{1}, bb name={battery}] (dp) {};
%       \node [black, left = 0.2 of dp] {$\cCapacity$};
%       \node [black, right = 0.2 of dp] {$\cMass$};
%   \end{tikzpicture}
%   \]
%   The concrete representation of the design problem is
%   \[
%       \text{battery} \colon\cCapacity\op\times \cMass \toinPos \Bool,
%   \]
%%
%   In this case, the feasibility question is
%   \[
%           \text{battery}(c^*, m) = \true \quad \Leftrightarrow  \quad
%            \text{a battery of mass $m$ is sufficient to provide the capacity $c^*$.}
%   \]
%   It is easy to convince oneself that this ``$\text{battery}$'' is monotone from basic physics consideration.
%   Monotonicity is equivalent to the following two assertions:
%%
%   \newcommand{\qsmall}[1]{{\color{blue}#1}}
%   \newcommand{\qlarge}[1]{{\color{red}#1}}
%%
%\begin{enumerate}
%\item We can provide a smaller capacity with the same mass:
%   \begin{eqnarray}
%       \text{For all}\  \qlarge{c_2}  \geq_{\cCapacity} \qsmall{c_1},\quad
%       \text{a battery of mass $m$ is sufficient to provide the capacity $\qlarge{c_2}$} \\
%       \Rightarrow
%       \text{a battery of mass $m$ is sufficient to provide the capacity $\qsmall{c_1}$.}
%   \end{eqnarray}
%%
%\item A battery of larger mass can provide the same capacity:
%   \begin{eqnarray}
%       \text{For all}\ \qlarge{m_2} \geq_{\cMass} \qsmall{m_1},\quad
%       \text{a battery of mass $\qsmall{m_1}$ is sufficient to provide the capacity $c^*$} \\
%       \Rightarrow
%       \text{a battery of mass $\qlarge{m_2}$ is sufficient to provide the capacity $c^*$.}
%   \end{eqnarray}
%\end{enumerate}
%
%   Assume that there is a linear relation between mass and capacity,
%   and such relation is described by the energy density~$\rho$ [Wh/kg].
%   Then the minimum mass to provide~$c^*$ is~$m_\text{min} = m_0 + c^* / \rho$.
%   So we have
%   \[
%       \text{battery}(c^*, m) = \true \quad \Leftrightarrow \quad  m \geq m_0 + c^* / \rho.
%   \]
%   A visualization of the feasible set $\feasibleset{\text{battery}}$ is in \cref{fig:battery-1}.
%   Monotonicity means that if we fix one point $\tup{c^*,m}$,
%   if we increase $c^*$, we can only see at most one transition, from feasible to unfeasible.
%   Similarly, if we increase $m$, there is at most one transition from unfeasible to feasible.
%
%\end{example}
%
%\begin{figure}[h!]
%    \caption{}
% \label{fig:battery-1}
%\end{figure}


\subsection{Querying}
\label{sec:dp-querying}

%\linkvideo{spring2021-functorial-comp-b:solving-queries} % Solving DP queries

\todotextjira{172}{\bernina: Explain briefly that in this section we show how the notion of querying introduced earlier connects to the mathematical definition of a \SY{design problem} as a \SY{boolean profunctor}.}


Sometimes we are not focused on whether a specific pair of resources, $\tup{\fun, \res}$ is a feasible pair for a given design problem (or monotone relation), but rather, we would like to know, for a fixed functionality $\fun$, what are \emph{all} the requirements $\res$ that render $\tup{\fun, \res}$ feasible. Or, analogously, we might want to know, for a fixed requirement $\res$, what are \emph{all} possible functionalities $\fun$ such that $\tup{\fun, \res}$ is feasible. We call this type of a question a \emph{query}. 


Equation~\cref{eq:dpi_to_dp_1} on the one hand, and \cref{eq:dpi_to_dp_2} on the other hand, give two perspectives on the mathematical definition of what we are calling a \emph{design problem}.
These two perspectives are analogous to something we already discussed in \cref{rem:rel-three-fun-descriptions}, when talking about binary relations.
There, we said that a binary relation from a set~\setA to a set~\setB is a subset~$\relA \setsubseteq \setA \cartprod \setB$,
but that such a relation~\relA can also, equivalently, be viewed as a function~$\phi_{\relA}\colon \setA \cartprod\setB \to \makeset{ \false, \true }$.
The subset~$\relA \setsubseteq \setA \cartprod \setB$ corresponded to the set
\begin{equation}
    \makeset{ \tupp{\ela, \elb} \setin \setA \cartprod \setB \mid \phi_{\relA}(\ela, \elb) = \true }.
\end{equation}
%
To make the analogy with \cref{eq:design-prob-as-upper-set} more precise, note that~\setA,~\setB,~$\makeset{ \false, \true }$, and~$\phi_{\relA}\colon \setA \cartprod\setB \sto \makeset{ \false, \true }$ live in the category of \emph{sets}, and that~$\funsp$,~$\ressp$,~$\Bool$ and~$\adp\colon \funspop \Ptimes \ressp\toinPos \Bool$ live in the category of \emph{posets}.

In \cref{rem:rel-three-fun-descriptions}, we also discussed two further ways to describe a relation~$\relA \setsubseteq \setA \cartprod \setB$: namely, we can transform the function~$\phi_{\relA}\colon \setA \cartprod\setB \sto \makeset{ \false, \true }$ either into a function
\begin{equation}
    \label{eq:curried-relation-again}
    \defmapperiod{
        \hat \phi_{\relA}}{\setA}{\sto
    }{
        \powerset(\setB)
    }{
        \ela
    }{
        \makeset{ \elb \setin \setB \mid \inrel{\ela
            }{
                \relA
            }{
                \elb} }
    }
\end{equation}
or a function
\begin{equation}
    \label{eq:co-curried-relation-again}
    \defmapperiod{
        \check \phi_{\relA}}{\setB}{\sto
    }{
        \powerset (\setA)
    }{
        \elb
    }{
        \makeset{ \ela \setin \setA \mid \inrel{\ela
            }{
                \relA
            }{
                \elb}}
    }
\end{equation}
%
There are analogous transformations for a \SY{design problem}~$\adp\colon \funspop \Ptimes \ressp\toinPos \Bool$.
Can you guess what they would be?

In order to use our ``sets to posets'' analogy and find an answer, it is useful to express the constructions we used in the setting of sets and relations entirely in terms of constructions from the category of sets, if possible.
Then the strategy is to identify what are the analogous constructions in the category of \SY{posets}, and this will allow us to make analogous definitions for \SY{design problems}.

The functions~$\hat \phi_{\relA}$ and~$\check \phi_{\relA}$ above have \SY{powersets} as their target objects.
What is the analogue of the \SY{powerset} operation in the category of posets?

The answer that we will use goes like this.
Given a set~\setA, there is a 1-to-1 correspondence between subsets of~\setA and functions~$\setA \sto \makeset{ \false, \true }$ (similar to above, a set corresponds here to its indicator function).
Thus, $\powerset (\setA)$ can be seen to correspond to~$\HomSet{\Set}{\setA}{\makeset{ \false, \true}}$.
The latter is definitely an expression we can transfer, by analogy, to the category of \SY{posets}, namely we can consider~$\HomSet{\Pos}{\posA}{\Bool}$.
And from \cref{ex:UpperSetsViaFunctors} we know that \SY{monotone maps}.
$\posA \mto \Bool$ correspond to \emph{upper} subsets of~\posA.
So~$\HomSet{\Pos}{\posA}{\Bool}$ corresponds to set~$\uppersets (\posA)$ of upper subsets of~\posA (c.f.
\cref{sec:UpperLowerSets} for the definitions of upper and \SY{lower sets}).

We now can write down the ``poset'' analogues of the functions~$\hat \phi_{\relA}$ and~$\check \phi_{\relA}$.
Namely, given a \SY{design problem} \cref{eq:generic-design-prob}, we have associated functions
\begin{equation}
    \label{eq:dp-left-curry}
    \widehat{\adp} \colon \funspop \mto \uppersets  (\ressp)
\end{equation}
and
\begin{equation}
    \label{eq:dp-right-curry}
    \widecheck{\adp} \colon \ressp \mto \uppersets (\funspop).
\end{equation}
However, we are not quite finished: are these \SY{monotone functions}?
Which \SY{poset} structure can we cho3ose on~$\uppersets (\ressp)$ and~$\uppersets (\funspop)$, respectively, so that~$\widehat{\adp}$ and~$\widecheck{\adp}$ are monotone?

\begin{gradedexercise}[\exname{CurryingDesignProblems}]
    \label{ex:CurryingDesignProblems}
    Let~$\adp\colon \funspop \Ptimes \ressp\toinPos \Bool$ be a design problem.
    In this exercise we will show that \cref{eq:dp-left-curry} corresponds to a monotone function
    \begin{equation}
        \funsp \mto \tupp{\uppersets (\ressp), {{\setsupseteq}}},
    \end{equation}
    and that \cref{eq:dp-right-curry} corresponds to a monotone function
    \begin{equation}
        \ressp \mto \tupp{\setOfLowersets(\funsp), \setsubseteq}.
    \end{equation}
    Here~$\uppersets (\ressp)$ denotes the set of \SY{upper sets} of~$\ressp$ and~$\setOfLowersets(\funsp)$ denotes the set of \SY{lower sets} of~$\funsp$.

    \begin{enumerate}
        \item Show that~$\widehat{\adp} \colon \funspop \mto \uppersets (\ressp)$ and~$\widecheck{\adp} \colon \ressp \mto \uppersets (\funspop)$ are \SY{monotone maps} when we consider~$\uppersets (\ressp)$ and~$\uppersets (\funspop)$ to have the partial order corresponding to the inclusion of subsets.
        \item Show that the \SY{poset} $\tupp{\uppersets\funspop,\setsubseteq}$ and the \SY{poset} $\tupp{\setOfLowersets\funsp, \setsubseteq}$ are isomorphic.
        \item Show that there is a 1-to-1 correspondence between  \SY{monotone functions} \begin{equation}
                  \funspop \mto \tupp{\uppersets \ressp, \setsubseteq}
              \end{equation}
              and  \SY{monotone functions} \begin{equation}
                  \funsp \mto \tupp{\uppersets \ressp, {{\setsupseteq}}},
              \end{equation}
              where in the latter poset, the order is given by the relation of ``containment'' (as opposed to ``inclusion'').
        \item Explain, in a few words, why the above steps prove the stated goal of this exercise.
    \end{enumerate}
\end{gradedexercise}
\solutionof{CurryingDesignProblems}

\todotextjira{173}{\bernina: Explain how the above discussion connect to querying and intuition of the associated optimization problems.}

\subsection{A note on pre-orders}
    The theory of \SY{design problems} can be easily generalized to pre-orders.
    This means that there could be two elements~$\posAel$ and~$\posBel$ such that~$\posAel\posAleq \posBel$ and~$\posAel \posAgeq \posBel$ but~$\posAel \neq \posBel$.

    This is actually common in practice.
    For example, if the order relation comes from human judgement, such as customer preference, all bets are off regarding the consistency of the relation.
    We will only refer to \SY{posets} for two reasons:
    \begin{enumerate}
        \item The exposition is smoother.
        \item Given a \SY{pre-order}, computation will always involve passing to the \SY{poset} representation.
    \end{enumerate}
    This means that, given a \SY{pre-order}, we can consider the \SY{poset} of its isomorphism classes, by means of the following \SY{equivalence relation}:
    \begin{equation}
        \posAel \simeq \posBel \quad \equiv \quad (\posAel \posAleq \posBel) \wedge (\posBel \posAleq \posAel).
    \end{equation}

